\chapter{Two-Level Cascading: Architecture, Calibration, and Theory}
\label{chap:cascading}

This chapter is the technical core of the report. It develops, precisely
and end to end, the two-level cascade that is layered on top of each
baseline to produce SUQL~V1 and Trummer~V1: the architecture
(Section~\ref{sec:cascade-architecture}), the per-variant instantiation
with the real prompts and parameters used by the implementation
(Section~\ref{sec:variant-instantiation}), the calibration procedure
that chooses the escalation threshold(s) from data
(Section~\ref{sec:beta-posterior}), the operator-reordering rule that
decides in which sequence the cheap and expensive stages should run
(Section~\ref{sec:operator-ordering}), and -- new in this revision -- a
formal proof that a calibrated cascade's expected and realized cost
(both LLM calls and wall-clock time) is lower than an all-expensive
baseline's once the candidate pool is large enough
(Section~\ref{sec:theory-cost}).

\section{Architecture}
\label{sec:cascade-architecture}

The cascade treats the semantic predicate $\sigma_U$ not as a single
opaque call, but as a two-stage pipeline over each \emph{unit of work}
$u$ (a row, for SUQL~V1; a row scored inside a batched pass, for
Trummer~V1):

\begin{enumerate}
  \item \textbf{Stage 1 (cheap).} A cheap scorer $f_1(u, \sigma_U) \to p
        \in (0,1)$ (or, for Trummer~V1, a direct three-way label) estimates
        whether $u$ satisfies $\sigma_U$, using a smaller/faster model
        than Stage~2.
  \item \textbf{Stage 2 (expensive).} A single-row \answerfn{}-style
        call, invoked \emph{only} on the subset of units whose Stage~1
        output falls inside a calibrated ambiguity band.
\end{enumerate}

\begin{figure}[h]
\centering
\resizebox{\textwidth}{!}{%
\begin{tikzpicture}[
  node distance=5mm and 5mm,
  box/.style={draw, rounded corners, align=center, minimum height=8mm, font=\scriptsize, fill=blue!6},
  cheap/.style={draw, rounded corners, align=center, minimum height=8mm, font=\scriptsize, fill=green!12},
  exp/.style={draw, rounded corners, align=center, minimum height=8mm, font=\scriptsize, fill=orange!14},
  dec/.style={draw, diamond, align=center, font=\scriptsize, fill=yellow!15, inner sep=1pt},
  arr/.style={-{Latex[length=2mm]}}
]
\node[box] (in) {Structurally\\filtered units\\$u_1,\dots,u_n$};
\node[cheap, right=of in] (cheap) {Stage 1: cheap\\scorer $\to$ log-odds\\$\ell(u_i)$};
\node[dec, right=14mm of cheap] (dec) {$\ell(u_i)$ in\\ambiguity band\\$[\tau_-,\tau_+]$?};
\node[exp, above right=6mm and 10mm of dec] (exp) {Stage 2: expensive\\operator (row-wise\\\answerfn{} or\\batched verify)};
\node[box, below right=6mm and 10mm of dec] (accept) {Accept /\\reject directly};
\node[box, right=22mm of dec] (out) {Merged\\result set};
\draw[arr] (in)--(cheap)--(dec);
\draw[arr] (dec) -- node[above,font=\tiny]{escalate} (exp);
\draw[arr] (dec) -- node[below,font=\tiny]{resolved} (accept);
\draw[arr] (exp) -| (out);
\draw[arr] (accept) -| (out);
\end{tikzpicture}}
\caption{Two-level cascade shared by SUQL~V1 (unit = row, $B=1$) and
Trummer~V1 (unit = row, scored inside a batched pass of $B=8$ rows,
escalated individually).}
\label{fig:cascade}
\end{figure}

\begin{algorithm}[h]
\caption{Two-level cascaded semantic filter}
\label{alg:cascade}
\begin{algorithmic}[1]
\Require rows $R_{\sigma_S}$ (post structured pushdown), predicate
$\sigma_U$, cheap scorer $f_1$, expensive operator $f_2$, calibration set
$C$, target quality $Q^\ast$
\State $(\tau_-, \tau_+) \gets \textsc{CalibrateBand}(f_1, C, Q^\ast)$
       \Comment{Section~\ref{sec:beta-posterior}}
\State $\mathit{Accept}, \mathit{Reject}, \mathit{Escalate} \gets \emptyset,\emptyset,\emptyset$
\For{$u \in R_{\sigma_S}$}
  \State $\ell \gets \operatorname{logit}(f_1(u, \sigma_U))$
  \If{$\ell \ge \tau_+$}
    \State $\mathit{Accept} \gets \mathit{Accept} \cup \{u\}$
  \ElsIf{$\ell \le \tau_-$}
    \State $\mathit{Reject} \gets \mathit{Reject} \cup \{u\}$
  \Else
    \State $\mathit{Escalate} \gets \mathit{Escalate} \cup \{u\}$
  \EndIf
\EndFor
\For{$u \in \mathit{Escalate}$} \Comment{Only ambiguous units reach Stage 2}
  \If{$f_2(u, \sigma_U) = \textsc{Yes}$}
    \State $\mathit{Accept} \gets \mathit{Accept} \cup \{u\}$
  \EndIf
\EndFor
\State \Return $\mathit{Accept}$
\end{algorithmic}
\end{algorithm}

\section{Per-Variant Instantiation}
\label{sec:variant-instantiation}

\subsection{SUQL~V1: cascaded predicate pushdown}
\label{sec:suql-v1}
After structural pushdown, \texttt{project SUQL/v1/scorer.py} scores
every surviving row with a \emph{single-token, logprob-based} cheap
probe: the model is asked to emit exactly one token (\texttt{1} for Yes,
\texttt{0} for No), with \texttt{max\_tokens=1}, and the log-odds score
is extracted directly from the token-level log-probabilities of that one
completion token (the difference of the log-probabilities assigned to
the \texttt{1} and \texttt{0} tokens), rather than from a model-reported
confidence string. This is a deliberately more principled design than
asking the model to self-report a probability: token log-probabilities
are the model's own calibrated (if imperfect) belief, not a number it
has to compose in text.

\begin{lstlisting}[style=prompt, caption={Cheap Stage~1 scoring prompt, verbatim from \texttt{project SUQL/v1/scorer.py}.}, label={lst:suql-v1-cheap}]
Act as a high-recall first-pass semantic filter. Decide
whether this review contains review-specific evidence for a
Yes answer. Count direct wording, synonyms, described
examples, and reasonable implications as evidence. Give
borderline but condition-specific evidence the benefit of
the doubt. Do not discard an explicit mention because it is
negated, qualified, quoted, or critical: this is evidence
retrieval, not sentiment scoring. Do not count genre/topic
alone, generic praise or criticism, or mere lack of
contradiction as evidence.
Return exactly one token: 1 for Yes, 0 for No.

Question: {question}
Review: {review[:1800]}

Answer:
\end{lstlisting}

Rows with log-odds outside the calibrated ambiguity band
$[\tau_-,\tau_+]$ are resolved without any further LLM call; only rows
inside the band are escalated to a full expensive-stage call
(\texttt{cascade\_filter.py}'s \texttt{ANSWER\_SYSTEM}, essentially the
same recall-first prompt as the SUQL baseline's, Listing~\ref{lst:suql-prompt}).
Because the cheap stage here scores \emph{one row per call} ($B=1$, no
batching), it is strictly a \emph{subset} of the row-wise baseline's
expensive-stage calls that get replaced -- for a fixed band, SUQL~V1 can
only use fewer or equal \emph{expensive} calls than the SUQL baseline,
never more. Empirically this holds (Chapter~\ref{chap:experiments}: 
$30.05$ expensive calls in the cascade vs.\ effectively $|R_{\sigma_S}|=40$
in the baseline), but wall-clock time still grows substantially: every
row now also pays for one \emph{cheap}-stage call, and, as
Section~\ref{sec:theory-cost} proves from first principles, that
cheap-stage call is \emph{not actually cheap in wall-clock time} in the
current deployment -- its measured latency exceeds the expensive
stage's own per-call latency, because a single-token completion still
pays the full prompt-\emph{prefill} cost of a review-length input, and
prefill, not generation, dominates latency for a 1-token completion.
This is the single most important empirical/theoretical finding of this
chapter, developed fully in Section~\ref{sec:theory-cost}.

\subsection{Trummer~V1: cascaded, structured, batched join}
\label{sec:trummer-v1}
Trummer~V1 (\texttt{project Trummer/v1/trummer\_join/}) restructures the
pipeline relative to the Trummer baseline in three ways:
\begin{enumerate}
  \item An explicit structured filter (\texttt{structured\_filter.py})
        is executed \emph{before} any LLM call -- the correction of
        Section~\ref{sec:trummer-baseline} applies here too: this is the
        one variant in the Trummer family that actually gets
        SUQL-style pushdown.
  \item Surviving candidates are batched (\texttt{cheap\_batch\_size=8})
        and scored with a single cheap prompt per batch of 8 rows,
        asking for a direct three-way label per row --
        \texttt{YES}/\texttt{NO}/\texttt{UNCERTAIN} -- rather than a
        continuous probability.
  \item Rows whose cheap-stage \texttt{UNCERTAIN} label or whose
        $|\text{score}|$ falls below a learned confidence threshold are
        re-verified in an expensive batched pass
        (\texttt{expensive\_batch\_size=32}).
\end{enumerate}

\begin{lstlisting}[style=prompt, caption={Cheap Stage~1 batched prompt, verbatim structure from \texttt{project Trummer/v1/trummer\_join/cascade.py} (\texttt{\_cheap\_batch\_prompt}).}, label={lst:trummer-v1-cheap}]
Classify every explicit candidate pair below.
Predicate: {predicate}
{CHEAP_EVIDENCE_INSTRUCTIONS}
Return one line per candidate as CANDIDATE_ID: YES, NO, or
UNCERTAIN.
Do not include confidence, evidence, explanations, or prose.
Return every candidate exactly once and do not explain.

CANDIDATE_1:
Movie: {movie_1_text}
Review: {review_1_text}

CANDIDATE_2:
Movie: {movie_2_text}
Review: {review_2_text}
...
Decisions:
\end{lstlisting}

\begin{lstlisting}[style=prompt, caption={Expensive Stage~2 batched verification prompt, verbatim structure from \texttt{cascade.py} (\texttt{\_expensive\_batch\_prompt}).}, label={lst:trummer-v1-expensive}]
Evaluate the explicit candidate pairs below.
Predicate: {predicate}
{EXPENSIVE_EVIDENCE_INSTRUCTIONS}
Return one line per candidate as PAIR_ID: YES or PAIR_ID: NO.
Do not include confidence, evidence, explanations, or prose.
Return every candidate exactly once. Do not add prose, JSON,
or markdown.

PAIR_1:
Movie: {movie_1_text}
Review: {review_1_text}
...
Decisions:
\end{lstlisting}

This is why Trummer~V1 is faster than the Trummer baseline (an explicit
structured filter shrinks the candidate pool before any LLM call runs
at all, and both cascade stages amortize their prompt overhead across
$B=8$ rows per call) \emph{and} more accurate (the expensive
verification pass repairs exactly the incomplete-coverage failure mode
that the join-based baseline suffers from) at the same time.

\textbf{An honest asymmetry.} SUQL~V1's threshold is calibrated with the
Beta-posterior credible-bound procedure of Section~\ref{sec:beta-posterior}
(\texttt{profiler.py}'s \texttt{beta\_lower\_bound}). Trummer~V1's threshold,
as currently implemented (\texttt{cascade.py}'s \texttt{\_learn\_threshold}),
instead performs a simpler \emph{point-estimate agreement search}: it
scans candidate thresholds and returns the smallest one whose calibration-set
agreement with the expensive oracle meets the target, with no Beta-posterior
credible interval around that agreement estimate. Both are legitimate
calibration procedures -- and the general cost theorem of
Section~\ref{sec:theory-cost} applies to either, since it only requires
that the calibration procedure fix \emph{some} per-row escalation
probability $\pi$ from held-out data -- but only the SUQL~V1 procedure
currently carries the finite-sample, confidence-level guarantee proven in
Proposition~\ref{prop:beta-bound}. Porting the Beta-posterior bound to
Trummer~V1's threshold search is flagged as near-term future work
(Chapter~\ref{chap:conclusion}).

\section{Calibrating the Cascade}
\label{sec:beta-posterior}

\subsection{Why not a hand-picked threshold?}
A na\"ive cascade would escalate everything with $|\ell_i| < \epsilon$
for a hand-picked constant $\epsilon$. This gives no quantitative
guarantee about the quality of the decisions made \emph{outside} the
band. We instead choose the band from a held-out \emph{calibration set}
$C$ of $(u, \sigma_U)$ pairs with known ground truth -- in the current
implementation, $|C|=20$ (\texttt{calibration\_budget}), sampled from the
structurally-filtered candidates and labeled by the expensive/oracle
operator itself.

\subsection{Setting up the posterior}
Fix a candidate threshold $\tau$ and define the accept set at that
threshold on the calibration set,
\[
  A_\tau = \{\, u \in C : \ell(u) \ge \tau \,\}.
\]
Let $\mathit{TP}(\tau)$, $\mathit{FP}(\tau)$, $\mathit{FN}(\tau)$,
$\mathit{TN}(\tau)$ be the usual confusion-matrix counts of $A_\tau$
(and its complement) against the ground truth on $C$. We model whether
an individual accepted (resp.\ rejected) item is truly positive as an
i.i.d.\ Bernoulli trial with unknown success probability equal to the
\emph{true} precision (resp.\ recall) of threshold $\tau$, and place a
non-informative uniform prior, $\mathrm{Beta}(1,1)$, on that unknown
probability. Because the Beta distribution is the conjugate prior of
the Bernoulli/Binomial likelihood, the posterior after observing
$\mathit{TP}(\tau)$ successes and $\mathit{FN}(\tau)$ failures (for
recall) is again Beta, in closed form:
\[
  \pi\bigl(r \mid \mathit{TP}, \mathit{FN}\bigr)
  \;=\; \mathrm{Beta}\bigl(1 + \mathit{TP}(\tau),\; 1 + \mathit{FN}(\tau)\bigr),
\]
and symmetrically for precision, using $\mathrm{Beta}(1+\mathit{TP}(\tau),
1+\mathit{FP}(\tau))$.

\subsection{The credible bound, and a subtle quantile pitfall}
\label{sec:quantile-pitfall}
We want a \emph{pessimistic}, high-confidence lower bound on the true
recall achieved by threshold $\tau$: a number $\widehat{r}_\alpha(\tau)$
such that we can assert, at confidence level $\alpha$ (\eg $\alpha=0.95$),
that the true recall of $\tau$ is at least $\widehat{r}_\alpha(\tau)$.

\begin{proposition}[Beta-posterior credible lower bound]
\label{prop:beta-bound}
Let $r\mid \mathit{TP},\mathit{FN} \sim \mathrm{Beta}(1+\mathit{TP},1+\mathit{FN})$
be the posterior above. For a confidence level $\alpha \in (0,1)$,
\[
  \widehat{r}_{\alpha}(\tau) \;=\; I^{-1}_{1-\alpha}\bigl(1+\mathit{TP}(\tau),\, 1+\mathit{FN}(\tau)\bigr)
  \;=\; \texttt{betaincinv}\bigl(1+\mathit{TP}(\tau),\, 1+\mathit{FN}(\tau),\, 1-\alpha\bigr)
\]
satisfies $\Pr[r \ge \widehat r_\alpha(\tau) \mid \mathit{TP}, \mathit{FN}] = \alpha$,
where $I^{-1}_y(a,b)$ is the inverse regularized incomplete Beta function
(the value $x$ such that $I_x(a,b)=y$).
\end{proposition}
\begin{proof}
By definition of the quantile function, $x=I^{-1}_y(a,b)$ satisfies
$\Pr[r \le x] = I_x(a,b) = y$ for $r\sim\mathrm{Beta}(a,b)$. Taking
$y=1-\alpha$ gives $\Pr[r \le \widehat r_\alpha(\tau)] = 1-\alpha$, hence
$\Pr[r > \widehat r_\alpha(\tau)] = \alpha$, i.e.\ $\widehat r_\alpha(\tau)$
is exceeded with posterior probability $\alpha$ -- the one-sided lower
credible bound at level $\alpha$.
\end{proof}
We flag a pitfall caught only through unit testing against a synthetic
confusion matrix with a known closed-form bound: the implementation's
\texttt{alpha} parameter denotes the \emph{confidence level itself}
(\eg \texttt{alpha=0.95}), not the tail probability; the tail probability
to plug into the quantile function is therefore $1-\alpha$, not $\alpha$.
Using the raw $\alpha$-quantile under this convention would silently
\emph{under}-certify the achieved recall whenever the calibration sample
is small. The precision bound is obtained symmetrically,
\[
  \widehat{p}_{\alpha}(\tau) \;=\; \texttt{betaincinv}\bigl(1+\mathit{TP}(\tau),\, 1+\mathit{FP}(\tau),\, 1-\alpha\bigr).
\]

\subsection{Choosing the band}
\label{sec:choosing-band}
The escalation band $[\tau_-, \tau_+]$ is the largest interval such
that, \emph{outside} the band, both certified quantities exceed a
target quality floor $Q^\ast$ (the cascade target, \eg $0.9$):
\[
  \tau_+ = \min\{\tau : \widehat{p}_\alpha(\tau) \ge Q^\ast\},
  \qquad
  \tau_- = \max\{\tau : \widehat{r}_\alpha(\tau) \ge Q^\ast\}.
\]
Everything with $\ell(u) \notin [\tau_-, \tau_+]$ is therefore
\emph{provably} (at confidence $\alpha$, on the calibration distribution)
safe to resolve without the expensive operator; everything inside the
band is exactly the set for which the cheap stage cannot certify
quality, and is exactly the set that should be escalated.

\section{Operator Ordering}
\label{sec:operator-ordering}

When more than one semantic sub-operator is chained in sequence -- as in
Trummer~V1's structural-filter~$\to$~batch-scoring~$\to$~verification
pipeline -- the order in which they run affects total expected cost even
though it does not affect the final result set (each operator is a
selection, and selections commute). Model $k$ independent selection
operators $O_1, \dots, O_k$, each with a per-invocation cost $c_i$ and a
\emph{selectivity} $s_i \in (0,1)$ -- the probability that a unit passes
$O_i$ and therefore needs to be evaluated by the next operator in the
chain. If the operators are applied in short-circuit order (evaluation
of a unit stops the moment it is rejected by some $O_i$), the expected
cost per input unit of evaluating them in the order $O_{\pi(1)}, \dots,
O_{\pi(k)}$ is
\[
  \mathbb{E}[\text{cost}] \;=\; \sum_{j=1}^{k} c_{\pi(j)} \prod_{m<j} s_{\pi(m)}.
\]
This expression is minimized, over all permutations $\pi$, by sorting
the operators in \emph{ascending} order of $c_i/(1-s_i)$, \ie operators
that are both cheap and likely to \emph{reject} should be evaluated
first. This is the same reordering rule used by ELEET's cascade
optimizer~\cite{eleet} and by classical selectivity-ordered join/filter
optimization; it directly justifies Trummer~V1's pipeline order: the
free structural filter runs first (cheapest, and typically rejects most
of the pool for a narrow predicate), the batched cheap semantic pass
runs second, and the batched expensive verification runs last, since by
construction (Section~\ref{sec:choosing-band}) it is only ever reached
by the already-narrowed ambiguous residual.

\section{Theoretical Cost Analysis of Calibrated Cascading}
\label{sec:theory-cost}

This section proves, from an explicit per-row cost model, that a
calibrated two-level cascade has lower expected -- and, for large enough
candidate pools, lower \emph{realized}, with high probability -- LLM-call
and wall-clock cost than an all-expensive (``vanilla'') baseline
evaluated over the same candidate pool, provided the calibration
achieves a non-trivial escalation rate relative to the cheap stage's own
batching factor. We then instantiate the theorem with the measured
constants of Chapter~\ref{chap:experiments} and show it correctly
predicts -- not merely describes after the fact -- the qualitative
asymmetry between SUQL~V1 (which fails the theorem's margin condition,
and indeed loses on wall time) and Trummer~V1 (which satisfies it
comfortably, and wins on both axes).

\subsection{Setup}
Consider a query whose structured pushdown yields $n$ candidate rows
$r_1,\dots,r_n$, modeled as i.i.d.\ draws from a fixed row distribution
$\mathcal D$ conditioned on the structured predicate -- a standard
exchangeability assumption for a fixed predicate applied to a sample of
the corpus. Compare two strategies for the semantic predicate $\sigma_U$:

\begin{description}
\item[Vanilla (all-expensive).] Every row is evaluated once by the
  expensive operator, at empirical cost $c_e$ calls and $\tau_e$
  wall-clock seconds per row (serial execution).
\item[Calibrated two-level cascade.] Every row is first evaluated by a
  cheap operator that processes rows in batches of size $B\ge 1$, at
  cost $c_c$ calls and $\tau_c$ seconds \emph{per batch} (so $c_c/B$
  calls and $\tau_c/B$ seconds amortized per row). A threshold pair
  calibrated on an $m$-row calibration set (itself labeled by the
  expensive oracle, at one-time cost $K_{\mathrm{calls}}=m\,c_e/B_{\mathrm{cal}}$
  calls and $K_{\mathrm{time}}=m\,\tau_e/B_{\mathrm{cal}}$ seconds, where
  $B_{\mathrm{cal}}$ is the calibration oracle's own batching factor)
  escalates each row independently to the expensive operator with
  probability $\pi \in [0,1]$, the row's calibrated escalation
  probability.
\end{description}

Let $S_n = \sum_{i=1}^n \mathbb 1[\text{row } i \text{ escalated}]$; under
the i.i.d.\ assumption, $S_n \sim \mathrm{Binomial}(n,\pi)$. The two
strategies' costs, as random variables in $n$, are
\begin{align*}
  \mathrm{Calls}_{\mathrm{van}}(n) &= n\,c_e, &
  \mathrm{Time}_{\mathrm{van}}(n) &= n\,\tau_e, \\
  \mathrm{Calls}_{\mathrm{casc}}(n) &= K_{\mathrm{calls}} + n\frac{c_c}{B} + S_n\,c_e, &
  \mathrm{Time}_{\mathrm{casc}}(n) &= K_{\mathrm{time}} + n\frac{\tau_c}{B} + S_n\,\tau_e.
\end{align*}

\subsection{Expected-cost crossover}
\begin{proposition}[Finite crossover point]
\label{prop:crossover}
Define the per-row margins
\[
  \Delta_{\mathrm{calls}} := (1-\pi)c_e - \frac{c_c}{B}, \qquad
  \Delta_{\mathrm{time}} := (1-\pi)\tau_e - \frac{\tau_c}{B}.
\]
If $\Delta_{\mathrm{calls}} > 0$, then
\[
  \mathbb E[\mathrm{Calls}_{\mathrm{casc}}(n)] < \mathrm{Calls}_{\mathrm{van}}(n)
  \iff n > n^\ast_{\mathrm{calls}} := \frac{K_{\mathrm{calls}}}{\Delta_{\mathrm{calls}}},
\]
and the gap $\mathrm{Calls}_{\mathrm{van}}(n) - \mathbb E[\mathrm{Calls}_{\mathrm{casc}}(n)]
= n\Delta_{\mathrm{calls}} - K_{\mathrm{calls}}$ grows without bound as
$n\to\infty$. The symmetric statement holds for time with
$n^\ast_{\mathrm{time}} := K_{\mathrm{time}}/\Delta_{\mathrm{time}}$. If
$\Delta_{\mathrm{calls}} \le 0$ (resp.\ $\Delta_{\mathrm{time}}\le 0$), no
finite $n$ makes the cascade's \emph{expected} cost lower: the cascade is
asymptotically worse, regardless of pool size.
\end{proposition}
\begin{proof}
By linearity of expectation and $\mathbb E[S_n]=n\pi$,
$\mathbb E[\mathrm{Calls}_{\mathrm{casc}}(n)] = K_{\mathrm{calls}} + n(c_c/B + \pi c_e)$.
The inequality
$K_{\mathrm{calls}} + n(c_c/B+\pi c_e) < n c_e$
rearranges to $n\bigl[(1-\pi)c_e - c_c/B\bigr] > K_{\mathrm{calls}}$, \ie
$n\Delta_{\mathrm{calls}} > K_{\mathrm{calls}}$, which (for
$\Delta_{\mathrm{calls}}>0$) holds iff $n > K_{\mathrm{calls}}/\Delta_{\mathrm{calls}}$;
for $\Delta_{\mathrm{calls}}\le 0$ the left side is non-increasing in
$n$ and the inequality never holds for $n$ large. The time statement is
identical with $(\tau_c,\tau_e,K_{\mathrm{time}})$ in place of
$(c_c,c_e,K_{\mathrm{calls}})$.
\end{proof}

\subsection{High-probability realized-cost dominance}
Proposition~\ref{prop:crossover} bounds the cascade's \emph{expected}
cost. Because $S_n$ is random, one might worry that an unlucky
realization -- an unusually ambiguous sample -- makes the cascade's
\emph{actual} cost exceed vanilla's even for $n>n^\ast$. Hoeffding's
inequality shows this risk vanishes exponentially fast in $n$.

\begin{proposition}[Exponentially vanishing failure probability]
\label{prop:hoeffding}
For $n > n^\ast_{\mathrm{calls}}$, let
$t_n := (n\Delta_{\mathrm{calls}} - K_{\mathrm{calls}})/c_e > 0$. Then
\[
  \Pr\bigl[\mathrm{Calls}_{\mathrm{casc}}(n) \ge \mathrm{Calls}_{\mathrm{van}}(n)\bigr]
  \;\le\; \exp\!\left(-\frac{2t_n^2}{n}\right)
  \;=\; \exp\!\left(-\frac{2\bigl(n\Delta_{\mathrm{calls}}-K_{\mathrm{calls}}\bigr)^2}{n\,c_e^2}\right),
\]
which is $O(e^{-cn})$ as $n\to\infty$ for a constant
$c = 2\Delta_{\mathrm{calls}}^2/c_e^2 > 0$. The symmetric statement holds
for time.
\end{proposition}
\begin{proof}
$\mathrm{Calls}_{\mathrm{casc}}(n) \ge \mathrm{Calls}_{\mathrm{van}}(n)$
iff $K_{\mathrm{calls}} + n c_c/B + S_n c_e \ge n c_e$ iff
$S_n \ge n\pi + t_n$ (substituting $t_n$ as defined and solving for
$S_n$). Since $S_n\sim\mathrm{Binomial}(n,\pi)$ is a sum of $n$ i.i.d.\
bounded $\{0,1\}$ variables with mean $n\pi$, Hoeffding's inequality
gives $\Pr[S_n \ge n\pi + t_n] \le \exp(-2t_n^2/n)$ for any $t_n>0$.
Substituting $t_n$ and expanding: for $n$ large,
$(n\Delta_{\mathrm{calls}}-K_{\mathrm{calls}})^2/n \sim \Delta_{\mathrm{calls}}^2\, n$,
linear in $n$, so the bound decays as $e^{-\Theta(n)}$.
\end{proof}
Propositions~\ref{prop:crossover}--\ref{prop:hoeffding} together give the
precise sense in which ``enough rows'' make cascading win: past a finite,
computable crossover $n^\ast$, the cascade wins \emph{in expectation}
(Proposition~\ref{prop:crossover}), and the probability that a single
realized run nonetheless loses to vanilla decays exponentially in the
number of rows past that point (Proposition~\ref{prop:hoeffding}) -- not
merely on average, but with a quantified, shrinking failure probability.

\subsection{Worked instantiation from the measured system}
\label{sec:worked-instantiation}
The margin conditions $\Delta_{\mathrm{calls}}>0$ and
$\Delta_{\mathrm{time}}>0$ are not automatic -- they depend on the
calibrated escalation rate $\pi$ \emph{and} the cheap stage's batching
factor $B$. Table~\ref{tab:theory-constants} instantiates both cascades
with the constants measured in the 10-query benchmark
(Chapter~\ref{chap:experiments}), using a ``lean vanilla'' reference
($c_e=1$ call/row; $\tau_e$ taken as the same cascade's own measured
expensive-call latency, since both stages use a comparable expensive
model class) so that the comparison isolates the cascade mechanism
itself from the baseline implementation's own non-essential per-row
overhead (\eg the SUQL baseline's \texttt{summary()} calls,
Section~\ref{sec:suql-baseline}).

\begin{table}[h]
\centering\small
\begin{tabular}{lrr}
\toprule
Quantity & SUQL~V1 & Trummer~V1 \\
\midrule
Batching factor $B$                     & $1$      & $8$ \\
Escalation rate $\pi$ (measured)         & $0.751$  & $0.125$ \\
Cheap cost per row $\tau_c/B$ (s)        & $20.858$ & $0.527$ \\
Expensive cost per call $\tau_e$ (s)     & $2.970$  & $2.795$ \\
Calibration overhead $K_{\mathrm{calls}}$ & $20$ (row-wise) & $1$ (batched) \\
\midrule
$\Delta_{\mathrm{calls}} = (1-\pi) - 1/B$      & $-0.751$ & $+0.750$ \\
$\Delta_{\mathrm{time}} = (1-\pi)\tau_e - \tau_c/B$ & $-20.12\text{ s}$ & $+1.919\text{ s}$ \\
Finite crossover $n^\ast_{\mathrm{calls}}$      & none (asymptotically worse) & $\approx 1.3$ rows \\
\bottomrule
\end{tabular}
\caption{Measured cascade constants (from \texttt{aggregate.csv},
Chapter~\ref{chap:experiments}) plugged into
Proposition~\ref{prop:crossover}. $c_c/B$ and $c_e$ are both $1$ call in
call-count units, so $\Delta_{\mathrm{calls}} = (1-\pi) - 1/B$ directly.}
\label{tab:theory-constants}
\end{table}

The two rows of Table~\ref{tab:theory-constants} tell essentially the
whole story of Chapter~\ref{chap:experiments} \emph{before any
experiment is run}:

\textbf{SUQL~V1 ($B=1$, $\pi=0.751$).} Both margins are negative. On
calls: with no batching, the cheap stage costs exactly $1$ call per row
-- the same as the lean vanilla reference -- so the cascade can only win
on calls if the escalation rate $\pi$ is below $1-1/B=0$, which is
impossible for $B=1$ unless $\pi=0$. On time, the cheap stage's own
per-row latency ($20.858$s) already exceeds the expensive stage's
per-call latency ($2.970$s), so even a hypothetical $\pi=0$ (never
escalate) would only break even on time asymptotically, and any positive
escalation rate makes it strictly worse. No finite $n$ can rescue this:
Proposition~\ref{prop:crossover} predicts the cascade is asymptotically
dominated on both axes, exactly matching the $7.99\times$ wall-time
regression and the fact that SUQL~V1's apparent call-count win in
Chapter~\ref{chap:experiments} is measured only against the (non-lean)
SUQL baseline, whose own per-row cost is inflated by ancillary
\texttt{summary()}/parsing overhead the lean model does not credit to
$c_e$.

\textbf{Trummer~V1 ($B=8$, $\pi=0.125$).} Both margins are comfortably
positive. Batching amortizes the cheap stage's fixed per-call latency
over $8$ rows, driving its effective per-row cost ($0.527$s) an order of
magnitude below the expensive stage's per-call latency ($2.795$s); the
calibrated escalation rate ($12.5\%$) is low enough that only a small
fraction of rows ever reach the expensive stage at all. The calibration
overhead is also batched ($K_{\mathrm{calls}}=1$, vs.\ $20$ for SUQL~V1's
row-wise calibration), giving a call-count crossover of
$n^\ast_{\mathrm{calls}} = 1/0.75 \approx 1.3$ rows -- the cascade is
expected to win on calls almost immediately, which is exactly what
Chapter~\ref{chap:experiments} observes even at the modest $n=40$
structured candidates used in the benchmark.

\textbf{The practical corollary.} Batching the cheap stage is not a
minor implementation optimization -- it is the load-bearing structural
choice that determines whether the margin conditions of
Proposition~\ref{prop:crossover} can hold \emph{at all}. A row-wise
($B=1$) cheap stage needs a per-row latency below the expensive stage's
\emph{own} per-call latency to have any chance of winning on time; a
batched ($B>1$) cheap stage divides that requirement by $B$, trading a
harder per-row latency target for an easier per-batch one. This is
exactly the lever identified as future work for SUQL~V1
(Chapter~\ref{chap:conclusion}): batching its currently row-wise cheap
probe is, by this analysis, the single change most likely to convert it
from a quality-only cascade into one that also wins on cost, the way
Trummer~V1 already does.
